\section{Distinct Routes}
\label{sec:cses-distinct-routes}

\problemstatement{Given a directed graph of $n$ nodes and $m$ edges, find
the maximum number of \emph{edge-disjoint} routes from node $1$ to node
$n$, and output each route. Every edge may be used by at most one route.}

\begin{patternbox}
The maximum number of edge-disjoint paths from $s$ to $t$ equals the
\textbf{maximum flow} when every edge has capacity $1$ (Menger's theorem).
After computing the flow, \textbf{decompose} it into that many paths by
following edges that carry flow.
\end{patternbox}

\subsection*{Max flow, then decompose into paths}

Give every edge capacity $1$ and compute the max flow from $1$ to $n$; its
value $f$ is the number of edge-disjoint routes. An original edge $u\to v$
carries flow iff its forward capacity dropped to $0$ (equivalently its
reverse now has capacity $1$). To extract the routes, walk from the source:
at each node follow one unused flow-carrying out-edge, mark it consumed, and
continue to the sink; repeat $f$ times. Because the flow is edge-disjoint,
these walks partition the flow edges into $f$ paths.

\begin{ideabox}[Flow Decomposition Yields the Paths]
A unit $s$--$t$ flow of value $f$ always splits into $f$ edge-disjoint
$s\to t$ paths (plus possibly some cycles, which unit capacities here avoid).
Greedily tracing flow edges from the source recovers each path exactly once.
\end{ideabox}

\subsection*{Algorithm}

\begin{enumerate}
  \item Add each edge with capacity $1$; run max flow $1\to n$; let $f$ be
        its value.
  \item Collect, per node, the out-edges that carry flow.
  \item Trace $f$ paths from $1$ to $n$, consuming each flow edge once; print
        the count and paths.
\end{enumerate}

\subsection*{Dry run}

\dryrun{$1\!\to\!2$, $2\!\to\!4$, $1\!\to\!3$, $3\!\to\!4$; $n=4$.}

\begin{itemize}
  \item Max flow $=2$: routes $1\to2\to4$ and $1\to3\to4$.
  \item Both routes are edge-disjoint.
\end{itemize}

\subsection*{C++ implementation}

\begin{lstlisting}
#include <bits/stdc++.h>
using namespace std;

struct Dinic {
    struct Edge { int to; long long cap; int rev; };
    vector<vector<Edge>> g;
    vector<int> level, it;
    Dinic(int n) : g(n), level(n), it(n) {}
    void addEdge(int u, int v, long long c) {
        g[u].push_back({v, c, (int)g[v].size()});
        g[v].push_back({u, 0, (int)g[u].size() - 1});
    }
    bool bfs(int s, int t) {
        fill(level.begin(), level.end(), -1);
        queue<int> q; level[s] = 0; q.push(s);
        while (!q.empty()) {
            int u = q.front(); q.pop();
            for (auto& e : g[u])
                if (e.cap > 0 && level[e.to] < 0) { level[e.to] = level[u] + 1; q.push(e.to); }
        }
        return level[t] >= 0;
    }
    long long dfs(int u, int t, long long f) {
        if (u == t) return f;
        for (int& i = it[u]; i < (int)g[u].size(); i++) {
            Edge& e = g[u][i];
            if (e.cap > 0 && level[u] < level[e.to]) {
                long long d = dfs(e.to, t, min(f, e.cap));
                if (d > 0) { e.cap -= d; g[e.to][e.rev].cap += d; return d; }
            }
        }
        return 0;
    }
    long long maxflow(int s, int t) {
        long long flow = 0;
        while (bfs(s, t)) {
            fill(it.begin(), it.end(), 0);
            long long f;
            while ((f = dfs(s, t, LLONG_MAX)) > 0) flow += f;
        }
        return flow;
    }
};

int main() {
    int n, m;
    cin >> n >> m;
    Dinic dinic(n + 1);
    for (int i = 0; i < m; i++) {
        int a, b; cin >> a >> b;
        dinic.addEdge(a, b, 1);
    }

    long long f = dinic.maxflow(1, n);
    cout << f << "\n";

    // flow-carrying out-edges: forward edge with cap now 0 (capacity was 1)
    vector<int> ptr(n + 1, 0);
    for (long long r = 0; r < f; r++) {
        vector<int> path;
        int u = 1;
        path.push_back(1);
        while (u != n) {
            auto& es = dinic.g[u];
            while (ptr[u] < (int)es.size() &&
                   !(es[ptr[u]].to != 0 && es[ptr[u]].cap == 0 &&
                     dinic.g[es[ptr[u]].to][es[ptr[u]].rev].cap > 0))
                ptr[u]++;
            auto& e = es[ptr[u]];
            dinic.g[e.to][e.rev].cap = 0;             // consume this flow edge
            u = e.to;
            path.push_back(u);
        }
        cout << path.size() << "\n";
        for (int v : path) cout << v << " ";
        cout << "\n";
    }
    return 0;
}
\end{lstlisting}

\begin{complexitybox}
\textbf{Time Complexity.} $O(V^2E)$ for the flow; decomposition is $O(V+E)$
across all paths. \textbf{Space Complexity.} $O(V+E)$.
\end{complexitybox}

\begin{takeawaybox}
The most edge-disjoint $s$--$t$ paths is unit-capacity max flow (Menger's
theorem), and flow decomposition recovers the actual routes. This closes the
flow chapter: throughput, cut, matching, and disjoint paths are one
algorithm applied to four networks.
\end{takeawaybox}
